\section{Getting started with PCPOP}
This secion is divided in four parts, starting with how to solve classic non-commutative polynomial optimization and then moving to partially-commutative polynomial optimization, then the graph product polynomial optimization and finally trace partailly-commutative polynomial optimization. 
\subsection{Non-commutative polynomial optimization}
Consider the following problem with hermitian variables $X_1,X_2$ from \cite{} and already solved using Ncpol2sdpa.
\begin{equation}
  \label{ncpop_pna}
\begin{split}
       p^* &=\underset{\underline{X} , \psi , \mathcal{H}}{inf} \: \braket{\psi, X_1 \mult X_2 + X_2 \mult X_1\psi}, \\
    & s.t \quad \qquad X_1^2 - X_1 = 0, \\
    & \qquad \quad -X_2^2+X_2+1/2 \geq 0, \\
    & \qquad \quad \braket{\psi , \psi} = 1. \\ 
\end{split}
\end{equation} 
The julia code to solve the above problem using PCPOP is given below.
\begin{code}
    # To define the non-commutative variables X1,X2 and its monoid M
    @ncmonoid M X1 X2
    #  To define the equality constraints
    Projector(X1)
    build(M)

    # To define the objective function, inequality constraints and level of relaxation
    obj=X1*X2+X2*X1
    op_ge=[X2-X2^2+0.5]
    level=2

    # To solve the non-commutative polynomial optimization problem
    ov,model,var_dict,principal_matrix=npa(obj,level;op_ge=op_ge)

    println("Optimal value is ",ov)
\end{code}
To solve a non-commutative polynomial optimization problem using PCPOP, we follow the follwoing steps 
\begin{itemize}
    \item First we need to define the non-commutative variables and an object that stores information about the corresponding monoid using the macro \texttt{@ncmonoid}. The first argument is always for the monoid object and the rest for the variables. 
    \item Then we define the equality constraints using \texttt{add\_relations!}. For constriants like \emph{projector,Unitary,Unipotent} we can use built-in functions \texttt{Projector,Unitary,Unipotent} respectively. 
    \item Next, we define the objective function, inequality constraints and level of relaxation.
    \item Finally, we call the function \texttt{npa} to solve the optimization problem. \texttt{npa} takes the objective polynomial and the level as compulsory arguments and also some optional arguments like arrays of operator inequalities(\texttt{op\_ge}), operator equalities(\texttt{op\_eq}), expectation value equalities(\texttt{tr\_eq}) and inequalities(\texttt{tr\_ge}). The function returns the optimal value, the SDP model, a dictionary that maps the monomial basis to coressponding JuMP variables and finally the principal matrix.
\end{itemize}
PCPOP provides nice interface to declare variables. We can create a container of variables using \texttt{X[n,m]} where $X$ is the name for variable array,  $n$ is the number of hermitian variables and $m$ is the number of non-hermitian variables. In that case $X$ stores the hermitian variables and $X\_$ stores the non-hermitian variables. 
\begin{lstlisting}[style=code, mathescape=true]
julia> $\text{@ncmonoid M X[2,1] a\_}$ 
julia> $\text{X}$
$\text{2-element Vector} \{cgb.Variable\}:$
$X_1$
$X_2$
julia> $\text{X\_}$
$\text{1-element Vector} \{cgb.Variable\}:$
$X_{1\_}$
julia>$\text{X\_[1]}'$
$X_{1\_\_}$
julia> $\text{cgb.is\_herm(a\_)}$
$\text{false}$
julia> $\text{a\_}'$
$a_{\_\_}$

\end{lstlisting}
We can add relations using the built-in function \texttt{add\_relations!} which takes as input, an array of non-commutative polynomials. When the monoid is build using \texttt{build(M)}, the \texttt{AbstractAlgebra} package is used to compute the Gr\"obner basis for the ideal generated by the relations and then it is used to reduce the polynomials. It is illustrated below.

\begin{lstlisting}[style=code, mathescape=true]
julia>$\text{@ncmonoid M a b}$
julia>$\text{cgb.add\_relations!}(M, [a^2 - a-1,a*b-b*a])$
julia>$\text{cgb.build}(M)$
julia>$\text{a*b*a*a}$
$2.0b*a + 1.0b$
\end{lstlisting}

\subsection{Partially-commutative polynomial optimization}